% Key references slide for the Department PINN Presentation.
% Build:  pdflatex references_slide.tex   ->  render to PNG via make_png.py
\documentclass[border=0pt]{standalone}
\usepackage[T1]{fontenc}
\usepackage{lmodern}
\usepackage{xcolor}
\usepackage{array}
\usepackage{amsmath}

% SA-big palette used across the project figures
\definecolor{sablue}{HTML}{1F77B4}
\definecolor{saorange}{HTML}{FF7F0E}
\definecolor{sagrey}{HTML}{5A5A5A}
\definecolor{sared}{HTML}{CC3333}

\newcommand{\refnum}[1]{\textcolor{sablue}{\bfseries #1}}
\newcommand{\reftitle}[1]{\textbf{#1}}
\newcommand{\refwhy}[1]{{\small\itshape\textcolor{sagrey}{#1}}}

\begin{document}
\begin{minipage}{22cm}
\setlength{\parskip}{0pt}
\centering
{\Huge\bfseries Key References}\\[2pt]
{\color{sablue}\rule{22cm}{2.2pt}}\\[10pt]

\raggedright
\renewcommand{\arraystretch}{1.35}
\begin{tabular}{@{}p{0.9cm}@{}p{20.6cm}@{}}

\refnum{1.} &
\reftitle{Raissi, Perdikaris \& Karniadakis (2019).} Physics-informed neural
networks: a deep learning framework for solving forward and inverse problems
involving nonlinear PDEs. \emph{J.\ Comput.\ Phys.} \textbf{378}, 686--707.\\
& \refwhy{Foundational PINN formulation --- the data + physics-residual loss used in every model here.}\\[4pt]

\refnum{2.} &
\reftitle{Yazdani, Lu, Raissi \& Karniadakis (2020).} Systems biology informed
deep learning for inferring parameters and hidden dynamics.
\emph{PLoS Comput.\ Biol.} \textbf{16}(11), e1007575.\\
& \refwhy{Closest template to this work: inverse PINNs on biological ODE networks; the local-sensitivity/Fisher route to identifiability.}\\[4pt]

\refnum{3.} &
\reftitle{Yang, Meng \& Karniadakis (2021).} B-PINNs: Bayesian physics-informed
neural networks for forward and inverse PDE problems with noisy data.
\emph{J.\ Comput.\ Phys.} \textbf{425}, 109913.\\
& \refwhy{Basis of the Bayesian inverse PINN: HMC over parameters gives posterior marginals instead of a point estimate.}\\[4pt]

\refnum{4.} &
\reftitle{Raue \emph{et al.} (2009).} Structural and practical identifiability
analysis of partially observed dynamical models by exploiting the profile
likelihood. \emph{Bioinformatics} \textbf{25}(15), 1923--1929.\\
& \refwhy{The identifiability verdict --- profile likelihood distinguishes the structural from the practical (high-WNT $\theta_P$) wall.}\\[4pt]

\refnum{5.} &
\reftitle{Wang, Yu \& Perdikaris (2022).} When and why PINNs fail to train:
a neural tangent kernel perspective. \emph{J.\ Comput.\ Phys.} \textbf{449}, 110768.\\
& \refwhy{Explains the loss-term imbalance and spectral bias behind the training pathologies; motivates adaptive loss weighting.}\\[4pt]

\refnum{6.} &
\reftitle{Tancik \emph{et al.} (2020).} Fourier features let networks learn high
frequency functions in low dimensional domains. \emph{NeurIPS} \textbf{33},
7537--7547.\\
& \refwhy{Fourier-feature time embedding --- required to capture the circadian RA forcing and the ATRA pulse.}\\[4pt]

\refnum{7.} &
\reftitle{Ji, Qiu, Zhang \& Deng (2021).} Stiff-PINN: physics-informed neural
network for stiff chemical kinetics. \emph{J.\ Phys.\ Chem.\ A} \textbf{125}(36),
8098--8106.\\
& \refwhy{Diagnoses PINN failure on stiff multiscale kinetics --- the same stiffness that shapes the WNT--RA--HOX residual.}\\

\end{tabular}
\end{minipage}
\end{document}
